\ifdefined\TriptychPrintEdition
\else
\section*{The Propers: Detailed Commentary}\label{per-element-sweep}
\fi
\sectionguard
\subsection*{Introit --- protection for a people desiring God's courts\enspace\properrefs{Int.}}\label{introit-treatment}

Psalm~83 joins longing, pilgrimage and protection. Its speaker loves the
Lord's tabernacles, faints for the courts, and sees the pilgrims advance toward
Sion. The prayer for the anointed face stands beside the preference for one
day in those courts. Protection serves communion with God: the worshipper
wants more than escape from danger. The whole psalm blesses dwelling and
praising, recalls the difficult journey, and ends with trust
(Ps.~83:2--13).

The Introit appoints Ps.~83:10--11a and then 2--3a. The nesting birds at the
altars, the heart and flesh rejoicing, and the preference for a low place in
God's house belong to the complete psalm, beyond the sung boundaries. The
words actually appointed give desire its object in the courts and its appeal
in the face of God's Christ. The title identifies the sons of Core; the body
names no particular earthly anointed king.

Augustine's petition at §13 asks that Christ become known to all. The face is
the manifestation of the Son, and the better day at §14 is everlasting life.
The Church's present longing amid tribulation stretches toward that day. His
house has God himself as its wall, and pilgrimage receives strength through
grace (\emph{Enarr. Ps.}~83 §§5--8, 12--15, Tweed's NPNF translation).
Theodoret instead identifies the face corporately with the saved people who
are Christ's body. He reads the better day through the benefit of attending
God's temple, applying the captives' desire to Christian worshippers
(PG~80:1543--1544). Augustine's eternal day and Theodoret's account of temple
attendance give distinct emphases to the same preference.

Hugh interprets the Son as a shield receiving wounds for sinners: the
Passion gives the protective image its christological force. His one day is
immutable eternity (\emph{Postilla}, Ps.~83:10--11, Ed1703 f.~221rb).
Bellarmine appeals to the Messiah's merits before the Father and reads the
preferred day as heavenly happiness, for which spiritual protection is
needed (O'Sullivan, p.~255). His appeal to merits differs from Augustine's
prayer that Christ be made known. At the entrance, these readings give the
lessons' renunciation a positive good: the Lord already desired is the Lord
whose kingdom must be sought first.

\subsection*{Collect --- continuing favor directs slipping mortality\enspace\properrefs{Coll.}}\label{collect-treatment}

\latin{Custódi} asks God to keep his Church through perpetual favor. The
reason follows: \latin{sine te lábitur humána mortálitas}. The singular
\latin{humána mortálitas} is the subject of slipping, of being withdrawn from
harms, and of being directed toward saving things. God's plural
\latin{auxíliis} supplies the continuing aid. Cummiskey's English makes
``grace'' the final clause's subject; the Latin directs mortal humanity
through those helps. The movement includes both removal from harm and a
positive saving destination.

The prayer's communal subject belongs beside the Epistle's social vices and
charity. Keeping the Church includes members whose hostility and faction
need correction. The same dependence accompanies the active commands to walk
and seek. Schuster understands the Collect through the necessity of divine
aid, then relates the Sunday's sacrifice and Communion to cleansing and
strengthening toward salvation (\emph{Sacramentary} III, pp.~136--138).

The three orations have a substantial common transmission. Wilson's
\emph{Gregorian Sacramentary under Charles the Great}, supplement XXXIII,
pp.~173--174, prints their complete text under Sunday~XV from the
ninth-century Vatican witnesses R (Reginae~337) and O (Ottobonianus~313).
Gros's editions corroborate their consecutive incipits in Gellone~II
(Montpellier ms.~18, nos.~1121--1123), Sant Ruf (nos.~427--429), and
Vilabertran (nos.~488--490), again under Sunday~XV. These later incipit
editions identify the set without supplying a complete collation of its
wording. Gellone~II is the manuscript Gros dates to 1000--1025, distinct from
the earlier Gelasian Gellone book, Paris BnF lat.~12048.

Wilson's marginal chant cues from O also show a different combination: the
Protector, Bonum and Inmittit group accompanies Sunday~XIII with
\latin{Panis quem ego}, while the Custodi trio at Sunday~XV has another chant
group. Later bracketed additions include Alleluia cues. Wilson favors Alcuin
as the supplement's compiler; Gros proposes Benedict of Aniane. Neither
attributes personal composition of these prayers. The checked supplement,
three incipit editions and Schuster commentary identify no composer or
founding date of the complete received Mass (Wilson, introduction
pp.~xvii--xxii, xxx, xliv--xlv; p.~173 notes; Gros, Gellone~II,
pp.~54--55, 75--76, 155).

\subsection*{Epistle --- the Spirit bears fruit through embodied freedom\enspace\properrefs{Ep.}}\label{epistle-treatment}

Galatians moves from Paul's defense of the Gospel through Abrahamic promise,
adoption and freedom to the conduct of Christ's people. Before the appointed
lesson, freedom is directed to mutual service through charity; after it,
Spirit-led persons restore the fallen gently and bear burdens. The works of
the flesh include religious disorder and social violence as well as sexual
and appetitive sins. Enmity, faction and envy tear apart the very community
called to serve (Gal.~1--6, especially 5:13--6:2).

Chrysostom argues that flesh here is depraved choice, not bodily substance.
The body serves in hearing, preaching and writing; strife and heresy arise
through moral agency. In explaining Gal.~5:17 he actually cites the
Introit's soul-longing verse, Ps.~83:3, to show that the soul desires.
Longing for God is not condemned with sinful appetite. The good action called
fruit requires human diligence and God's loving kindness; passions may still
trouble a person without ruling effectively (\emph{Commentary on Galatians},
chapter~5, at vv.~16--24, trans.~Alexander).

Augustine distinguishes undergoing contrary desire from carrying it to
completion. Christ's grace gives a stronger delight in justice through which
sinful consent is resisted; complete freedom from concupiscence belongs to
resurrection. Both lists are exemplary rather than exhaustive
(\emph{Epistolae ad Galatas expositio} §§47--51, PL~35:2139--2142,
Latin electronic presentations). Thus the continuing struggle in v.~17 and
the crucifixion of vices in v.~24 describe a life of real grace and real
resistance together.

Aquinas distinguishes fruit acquired as an end from fruit produced in action.
Here the Spirit produces delightful acts of virtue. Relative to future
beatitude those acts are like flowers; their present sweetness resembles
medicine whose pleasantness serves health. Charity begins the ordering of
the virtues, and their delight serves a further good. His next lecture gives
vigils, fasting and labor as bodily discipline while expressly forbidding the
destruction of bodily nature (\emph{Super Galatas} V, lect.~6 [87779],
lect.~7 [87780], Latin, Turin 1953 text).

The appointed Vulgate and Douay enumerate twelve fruit terms. Michael
W.~Holmes's SBLGNT v1.2 at Gal.~5:22--23 enumerates nine. Chrysostom's Greek reception has nine;
Aquinas and Anthony expound the received twelve-term Latin list. The checked
editions establish the difference but do not explain its manuscript history.
Paul's warning concerns those who practice the named vices, while the
Spirit's fruit gives the life the law does not condemn.

Anthony relates the twelve fruits to the cleansed Samaritan's grateful return:
unity, humility, poverty and acknowledgment of mercy shape fruitful
penitence. His distinction between continence and chastity belongs to that
Latin reception. The sermon labeled Fourteenth Sunday has Luke~17 and
\latin{Inclina Domine}, so Anthony is expounding this Epistle in a different
formulary (Sunday~XIV §§15--18, trans.~Spilsbury). The kingdom excluded by
Paul's warning is the kingdom positively sought in the Gospel and Communion.

\subsection*{Gradual --- confidence reaches beyond the helper\enspace\properrefs{Grad.}}\label{gradual-treatment}

The two comparisons in Ps.~117:8--9 stand between the Lord-as-helper
confessions and the account of encircling enemies. Israel, Aaron's house and
those who fear the Lord are summoned to thanksgiving; personal deliverance
opens toward the gates of justice and the Lord's house and altar. The
rejected stone and blessed arrival are later, unsung verses of that complete
context. The poem does not identify a particular historical procession or
speaker (Ps.~117:1--29).

Augustine includes angels among the princes: good human and angelic helpers
act through God, who made them good. Their ministry is received, while
ultimate trust belongs to its source (\emph{Enarr. Ps.}~117 §4).
Hugh hears the Church and martyrs in the psalm and similarly places the
helpers' power under God. His further moral reading concerns secular or
ecclesiastical patrons from whom people seek present glory: hope for the
crown exceeds trust for victory (\emph{Postilla}, Ps.~117:8--9,
units \{g\}--\{i\}).

Theodoret presses human mortality, changeable wills and limited power, even
where willingness exists; princely authority is brief (PG~80:1811--1812).
Bellarmine also begins from inability and unwillingness, then permits trust
in pious people and angels insofar as it refers to God acting through them
(O'Sullivan, p.~370). Theodoret's emphasis is human frailty; Augustine's is
God's action through good helpers. Neither makes every created assistance
unacceptable. The Collect's helps and the Offertory's protecting angel place
that dependence within the day's petitions and praise.

\subsection*{Alleluia --- joyful approach takes the shape of obedience\enspace\properrefs{All.}}\label{alleluia-treatment}

The appointed Ps.~94:1 is a communal call to rejoice in the Lord. The complete
psalm grounds that praise in creation, kingship and shepherding, then warns
against the wilderness generation's refusal to hear. Hebrews~3--4 expressly
receives the unappointed warning as an address to present faith and entry
into God's rest. Joy opens a psalm whose hearers must answer God's voice
(Ps.~94:1--11; Heb.~3:7--4:11).

Augustine asks how anyone approaches the God who is everywhere. Bad habits
estrange by unlikeness; forgiveness and renewal restore the divine image,
like a coin re-engraved. Jubilant praise expresses a joy beyond ordinary
words (\emph{Enarr. Ps.}~94 §§1--4). Hugh's moral application addresses the
Matins Invitatory: the singer inviting others must walk the commandments,
with attention and devotion instead of vanity (\emph{Postilla}, Ps.~94:1,
Ed1703 ff.~249rb--vb). Bellarmine unites heart and lips, drawing on Christ's
rejoicing and Mary's Magnificat; the 1866 page says the singers
``express their joy'' (O'Sullivan, p.~298).

Theodoret sets the psalm's dramatic voice among Josiah and the priests after
idolatry's overthrow, while expressly attributing prophetic composition to
David. He extends the triumphant opening especially to apostles and martyrs
(PG~80:1639--1640). That is his historical and spiritual interpretation,
not a dating of the Alleluia. The praise sung before the Gospel calls the
assembly to the faithful approach that its words require.

\subsection*{Gospel --- the kingdom orders work and bodily need\enspace\properrefs{Gosp.}}\label{gospel-treatment}

The treasure and eye sayings lead directly into the two masters. The
\latin{Ideo} that begins the teaching on solicitude joins anxiety to
allegiance: freedom from divided service makes room for trust in the Father.
Matthew's wider Sermon has already joined kingdom and bread in prayer and
will require the doing of the Father's will. Jesus teaches disciples with
crowds present in the Galilean narrative setting, on an unnamed mountain
(Matt.~4:23--5:2; 6:8--34; 7:21--29).

Augustine interprets the passage through one final end. Necessities are
sought for the kingdom rather than the kingdom for necessities. Chrysostom
distinguishes possessing wealth from serving it, with Job as steward rather
than slave. His account demands generosity, not admiration for accumulation
(Augustine, \emph{De sermone} II.15--17 §§49--58; Chrysostom,
\emph{Matthew} homily~21 §§1--4). God and mammon name incompatible mastery;
the saying does not abolish every subordinate human service.

The argument for providence begins with the greater gifts of life and body,
then proceeds to birds, stature, lilies and passing grass. Human worth,
anxiety's incapacity and the Father's knowledge converge in first seeking the
kingdom and justice. The Vulgate and Douay cubit concerns stature. Food,
drink and clothing are the antecedent of the things added, not an unlimited
inventory of desired possessions (Matt.~6:25--33).

Chrysostom's next homily carries the argument through the lilies and treats
the added necessities as lesser than the kingdom sought. Augustine's
examples include Christ's purse, collections, famine relief and work, then
Paul's hunger and nakedness. Responsible provision and actual deprivation
both occur within providence. The physician can give or withhold temporal
aids while directing the patient toward perpetual rest (Chrysostom,
homily~22 §§1--3; Augustine, II.17 §§57--58).

Anthony gives a direct moral reading of the lilies as penitents flourishing
amid worldly life: root, whiteness and fragrance become humility, chastity
and good reputation. He distinguishes necessary labor from divided care and
treats temporal loss as a possible trial. His opening reason/sensuality
interpretation is his moral application; Matthew names God and mammon. His
Fifteenth Sunday pairs this Gospel with Gal.~5:25 onward and
\latin{Miserere mihi}, another medieval assemblage
(Sunday~XV §§2--3, 7, 12--17, trans.~Spilsbury).

Francis de Sales offers related spiritual illumination when he commends
careful work without over-anxiety in \emph{Introduction to the Devout Life}
III.10. That chapter explicitly turns to Martha, Luke~10:41, rather than
expounding the appointed Matthew passage; III.14 begins with poverty of
spirit. The checked English presentation leaves its translator unidentified.
Luke~12:13--34 supplies a closer biblical parallel: its ravens and lilies
stand between the rich fool and almsgiving. Matthew's own context retains
necessary action and suffering, including persecution and the daily trouble
of unappointed v.~34.

\subsection*{Offertory --- protection surrounds those invited to taste\enspace\properrefs{Off.}}\label{offertory-treatment}

Psalm~33 moves from a poor man's deliverance to instruction in fear of the
Lord. After the protecting angel and the invitation to taste, it teaches
restraint of speech, departure from evil, doing good and seeking peace.
Its just person has many afflictions. The antiphon ends before v.~9's final
beatitude; neither that conclusion nor the lack-no-good promise of vv.~10--11
is sung here (Ps.~33:7--23).

Basil expressly gives each believer an attending angel whose protection
surrounds him; sinful conduct can drive the guardian away. His tasting
concerns experience of Christ the true bread, with the Lord's flesh as true
food. Honey's sweetness illustrates learning through experience, and a
present taste of grace awakens hunger for fuller enjoyment
(\emph{Homily on Psalm}~33 §§5--6, Greek and facing Latin,
PG~29:363--366).

Augustine identifies the angel as Christ, the messenger of the great
counsel, and explains tasting by receiving Christ's Body and Blood. His
continuation warns against measuring God's goodness by the wicked person's
prosperity (\emph{Enarr. Ps.}~33, second exposition §§10--13,
trans.~Tweed). His christological messenger differs from Basil's personal
guardian. Aquinas offers a ministering angel, Christ and a prelate guarding
the flock; he then explains inward tasting that precedes spiritual sight and
steadies affection (Ps.~33 nn.~8--9 [87112--87113]). That paragraph does
not explicitly name Eucharistic species.

Bellarmine reads a mighty created protector and personal experience of trust,
then nourishment after spiritual rebirth, with references to 1~Peter~2 and
Hebrews~6 (O'Sullivan, pp.~91--92). His nourishment paragraph likewise does
not explicitly identify Eucharistic species. First Peter itself reuses
tasting at 2:3 and the psalm's ethical teaching at 3:10--12. The invitation
at the offering therefore bears an explicit Eucharistic reception in
Augustine alongside the other witnesses' distinct readings; its spiritual
fruit is not measured by bodily sweetness.

\subsection*{Secret --- the saving offering cleanses sins\enspace\properrefs{Sec.}}\label{secret-treatment}

The Church asks that \latin{hæc hóstia salutáris} become both purification
of sins and propitiation before divine power. The two predicates concern the
saving offering. Cummiskey renders the first as an action cleansing the
worshippers and the second as rendering divine majesty propitious; the Latin
makes the offering \latin{purgátio} and \latin{propitiátio}. The petition
asks for divine saving action, with sin itself to be cleansed.

Schuster places this sacrifice within the Collect's dependence on divine help
and explicitly returns to it when explaining the Postcommunion
(\emph{Sacramentary} III, p.~138). The complete prayer stands between
\latin{Custodi} and \latin{Purificent} in Wilson's supplement XXXIII,
pp.~173--174; Gros's three incipit editions transmit the same sequence.
No personal composer of the prayer is identified in those witnesses.

The Epistle's crucifixion of vices and the Secret's purification meet at the
need for a changed life. Galatians itself names Christ's self-giving and
cross outside the appointed passage (Gal.~1:3--4; 2:19--21; 6:14--18).
The Mass's moral demands remain within the Church's request for a saving
offering. The requested cleansing is more than the worshipper's improved
composure; it concerns the sins from which the worshippers ask God to free
them.

\subsection*{Communion --- first seeking remains the rule of reception\enspace\properrefs{Comm.}}\label{communion-treatment}

\latin{Primum quǽrite regnum Dei} returns to the Gospel's last appointed
verse. The antiphon moves \latin{primum} to the front, omits \latin{ergo}
and \latin{et iustítiam eius}, reduces \latin{hæc ómnia} to \latin{ómnia},
and adds \latin{dicit Dóminus}. The complete Douay verse includes justice
and the demonstrative; it supplies the English Scripture witness beside the
Latin adaptation (Matt.~6:33; Missal no.~1580).

The Gospel still identifies the things added as food, drink and clothing.
At reception, its priority remains a command. The kingdom has not become a
means of obtaining the necessities; necessities serve life under God's
kingdom. Augustine's order of ends and Chrysostom's lesser added goods both
expound the full verse, as does Anthony's distinction between supreme and
secondary good (Augustine, II.16--17 §§55--58; Chrysostom, homily~22 §3;
Anthony, Sunday~XV §15). None comments on the Roman abbreviation as such.

The Epistle gives conduct its corresponding seriousness: charity and peace
belong to the Spirit's fruit, while practiced vice excludes from the kingdom.
The Communion's omission of the justice clause does not contradict the full
Gospel. The people who receive remain the people commanded to seek.

\subsection*{Postcommunion --- the sacraments purify, fortify and lead\enspace\properrefs{Postcomm.}}\label{postcommunion-treatment}

The final petition makes \latin{tua sacraménta} the subject of three verbs:
\latin{puríficent}, \latin{múniant}, \latin{ducant}. The recipients are
\latin{nos}, and the end is perpetual salvation. Cummiskey's historical
English says ``procure'' where the Latin asks that the sacraments lead. The
Church asks for continuing purification and strength toward that end; she
does not declare it already attained.

Schuster directly joins this prayer to the Secret, describing Communion's
strengthening of virtue and the mysteries' service to eternal salvation
(III, p.~138). The Catechism explains union with Christ, growth and cleansing,
charity, incorporation in the Church and commitment to the poor
(CCC~1391--1397). The poor belong explicitly to paragraph~1397. These
fruits concern the living bonds of Christ's people, just as the Epistle's
charity and peace answer enmity and faction.

The concluding prayer thus returns the Church to dependence after reception.
Her members have heard the kingdom's priority and Paul's demand for
Spirit-led action; they now ask to be purified, fortified and led. The
saving end remains a gift sought through perseverance, without a promise
that every recipient will feel consolation or possess temporal abundance.
